% Ameaças à validade — Lab03, Issue #49.
% Estrutura: 4 categorias de Wohlin et al. (2000) com evidências quantitativas do projeto.

\section{Ameaças à validade}

\subsection{Validade de construto}

\textbf{Release $\neq$ deploy.} A frequência de deploy e o \textit{lead time} são calculados a
partir de \textit{releases} publicadas no GitHub. Uma biblioteca que publica uma versão por mês
não ``faz deploy em produção''; quem faz isso são seus usuários. Repositórios de
infraestrutura ou de documentação podem não publicar nenhuma \textit{release} mesmo implantando
continuamente. Isso limita a generalização para projetos que adotam entrega contínua sem
\textit{releases} formais.

\textbf{Falha de CI $\neq$ falha em produção (CFR variante a).} O CFR(a) conta execuções de
CI com conclusão de falha; erros de lint, testes instáveis e falhas de ambiente não são
incidentes de produção. A mediana de 5,7\% pode subestimar (se falhas reais ficam ocultas)
ou superestimar (se a maioria das falhas não chega aos usuários).

\textbf{Heurística corretiva não validada (CFR variante b).} A v12 atingiu F1 = 0,689
(IC95\% 0,563--0,790) na amostra independente, abaixo da meta de 0,70. Toda conclusão
que use o CFR(b) — RQ03b, RQ05b e a combinação C2 da RQ07 — é exploratória. O $\kappa$
de Fleiss em \textit{corretiva} foi de 0,273 na amostra-ouro (razoável), o que indica que
a própria definição de ``release corretiva'' é ambígua para os avaliadores humanos.
As versões v8--v12 foram concebidas após examinar os erros da amostra de desenvolvimento,
tornando o F1 interno otimista.

\textbf{Rebase e squash distorcem a data do commit.} O \textit{lead time} usa
\texttt{commit.author.date} (E§3). Um rebase interativo altera a data do autor para o
momento do rebase, comprimindo artificialmente o LT para commits reorganizados. Não há
como detectar isso a partir da API pública.

\textbf{Releases fora do branch principal.} \textit{Releases} não são filtradas pelo
\textit{default branch}; backports publicados em ramos de manutenção entram no cálculo.
A contagem de \textit{releases} com \texttt{compare diverged} fica em
\texttt{n\_compare\_divergentes} no CSV.

\textbf{Correlação positiva DF × CFR(b) como artefato de denominador.} O coeficiente
$\rho = 0{,}62$ entre frequência de deploy e CFR(b) pode ser parcialmente explicado por
confundimento: projetos com muitas \textit{releases} têm mais oportunidades de publicar
correções, inflando o numerador do CFR(b) sem que isso reflita mais falhas reais.

\subsection{Validade interna}

\textbf{F1 otimista na escolha da heurística.} A versão escolhida (v12) foi selecionada
como a de maior F1 entre doze versões avaliadas na mesma amostra de desenvolvimento de
60 repositórios / 300 \textit{releases}. Sem um conjunto de validação separado, o F1
interno (0,722) é inflado; o F1 independente (0,689) é a estimativa mais confiável.

\textbf{Rotulagem assistida.} Os três avaliadores partilharam o protocolo de consenso e
discutiram desempates em reunião, o que reduz a independência real dos rótulos. O $\kappa$
de Fleiss por dimensão — 0,352 em ``entrega real'' e 0,273 em ``corretiva'' — indica
concordância razoável, mas aquém do que seria esperado em avaliadores verdadeiramente
independentes.

\textbf{Dados ausentes não aleatórios.} Os 23 repositórios sem tempo de recuperação
(nenhum episódio concluído) são projetos com poucas falhas de CI ou com episódios todos
censurados — possivelmente os mais estáveis da amostra. Excluí-los subestima o tempo
mediano de recuperação. Da mesma forma, os 2 repositórios sem CFR(b) são outliers de
comportamento incomum (\texttt{akto-api-security/akto}, com 1307 \textit{releases} avaliadas,
e \texttt{akveo/react-native-ui-kitten}, com 24 \textit{releases} todas censuradas).

\textbf{Censura no tempo de recuperação.} Episódios sem sucesso posterior até o fim da janela
são censurados e ficam fora da mediana. Se esses episódios são sistematicamente mais longos, a
mediana de 2,8~h subestima o tempo real de recuperação.

\subsection{Validade externa}

\textbf{Amostra por ordem alfabética com piso de estrelas.} Os 300 repositórios são os
primeiros que passam nos filtros (ações de CI, $\geq$ 5 \textit{releases}, $\geq$ 50 runs
válidos) em ordem alfabética dentro de faixas de estrelas com piso de aproximadamente 1000
estrelas. Isso exclui projetos menos populares e introduz viés de seleção: projetos com
nomes no início do alfabeto estão sobre-representados.

\textbf{Popularidade truncada.} O piso de popularidade garante que todos os repositórios
tenham alguma visibilidade, mas impede conclusões sobre projetos em estágio inicial ou de
nicho. Os resultados descrevem projetos open-source consolidados, não o ecossistema como um todo.

\textbf{Linguagens e tipos de projeto.} A amostra é dominada por TypeScript (22\%), Python
(18\%), Go (14\%) e Rust (10\%). Linguagens com ecosistemas de \textit{release} diferentes
(Java, PHP, Ruby) estão sub-representadas ou ausentes. O rótulo de tipo do projeto
(biblioteca, aplicação, CLI) está disponível apenas para os 60 repositórios da amostra-ouro.

\textbf{Janela de 12 meses.} A janela 2025-10-02 a 2026-10-02 captura um período específico;
projetos com ciclos de \textit{release} irregulares ou que mudaram de ritmo nesse intervalo
podem ter métricas não representativas do seu comportamento histórico.

\subsection{Validade de conclusão}

\textbf{Múltiplas comparações (RQ06).} A análise de fatores realiza 30 testes (5 fatores
$\times$ 6 métricas). A correção de Holm foi aplicada com $\alpha = 0{,}05$, e apenas 5
associações permaneceram significativas. Tamanhos de efeito medianos ($\eta^2_H$ entre 0,06
e 0,21) indicam que os fatores explicam pouca da variação total.

\textbf{Correlação de Spearman é monotônica, não linear.} O $\rho$ mede associação
monotônica; uma relação em U ou com limiar não seria detectada. A interpretação de
$\rho = 0{,}62$ entre DF e CFR(b) como ``correlação forte'' não implica causalidade nem
descarta formas funcionais não monotônicas.

\textbf{Número de testes por subgrupo.} Em RQ06, grupos com $n < 5$ foram excluídos.
O fator tipo do projeto tem apenas $n = 51$ repositórios com rótulo, contra $\approx$ 300
nos outros fatores, reduzindo o poder dos testes correspondentes.

\textbf{Sensibilidade à definição operacional (RQ07).} Entre 22\% e 40\% dos repositórios
mudam de categoria DORA ao trocar a combinação de métricas (C1, C2 ou C3), com $\kappa$
linear entre 0,46 e 0,72. Conclusões sobre a distribuição de níveis DORA dependem da
definição operacional escolhida, e comparações com outros estudos que adotem definições
diferentes devem ser feitas com cautela.
